\section*{Лекция 9. Разложение ошибки на смещение и разброс}

На практике регулярно происходит так, что одна конкретная модель (или класс моделей) показывает себя лучше других. В некоторых случаях бывает так, что линейная модель показывает себя лучше, чем решающие деревья, хотя мы недавно обсуждали разноплановость и эффективность дереьев. Но почему так происходит? Есть ли какое"=то обоснование тому, как подбираются модели для решения задачи?

В рамках данной лекции мы рассмотрим одну из наиболее важных тем, пожалуй, всего машинного обучения. Мы покажем, как можно интерпретировать ошибку модели на данных, и как мы можем контроллировать этот процесс.